\section{Proof of Proposition~\ref{prop:lbproduct}}\label{app:lbproduct}
\begin{lemma}[Isolation \cite{MulmuleyVaziraniVazirani1987}]\label{lem:isolation}
Let $\mathcal F$ be a nonempty family of subsets of a finite set $W$, and let
$w:W\to\{1,\dots,R\}$ be uniformly random. Then the minimum of
$w(S)=\sum_{e\in S}w(e)$ over $S\in\mathcal F$ is attained by a unique set with
probability at least $1-|W|/R$.
\end{lemma}

\begin{proof}
For $e\in W$ let $\alpha_e=\min\{w(S):S\in\mathcal F,e\notin S\}
-\min\{w(S\setminus\{e\}):S\in\mathcal F,e\in S\}$ (with
$\min\emptyset=+\infty$); $\alpha_e$ does not depend on $w(e)$. If two
distinct minimizing sets exist, some $e$ lies in one and not in the other,
and then $w(e)=\alpha_e$, an event of probability at most $1/R$.
\end{proof}

\begin{proof}
Take an instance of multicolored clique: $k$ color classes identified with
$\{1,\dots,N\}$ and, for $c<d$, edge lists
$R_{cd}=((a_1,b_1),\dots,(a_m,b_m))$ of distinct pairs. Choose independent
uniform weights $w(e)\in\{1,\dots,2|R|\}$ on the edges $e$ of
$R=\bigsqcup R_{cd}$, and let $W_0=1+\binom k2\,2|R|$. Use integer variables
$x_c\in[1,N]$ and, for each pair $c<d$ and $l=1,\dots,m=|R_{cd}|$,
$\zeta_l,S_l\in\{0,1\}$ and $A_l,B_l\in[0,N]$, with $S_0=A_0=B_0=0$. Let
\begin{align*}
\Phi&=\sum_{c<d}\Bigl[\sum_{l}\bigl((S_l-S_{l-1}-\zeta_l)^2
+(A_l-A_{l-1}-a_l\zeta_l)^2+(B_l-B_{l-1}-b_l\zeta_l)^2\bigr)\\
&\qquad\qquad+(S_m-1)^2+(A_m-x_c)^2+(B_m-x_d)^2\Bigr],\qquad
\Psi_w=W_0\Phi+\sum_{c<d}\sum_lw(a_l,b_l)\zeta_l .
\end{align*}
\emph{Decomposition.} Take a central bag $\{x_1,\dots,x_k\}$; for each pair a
bag $\{x_c,x_d,S_m,A_m,B_m\}$ adjacent to it, followed by the path of bags
$\{S_{l-1},A_{l-1},B_{l-1},\zeta_l,S_l,A_l,B_l\}$, $l=m,\dots,1$. Every term
lies in a bag, running intersection holds, and $p=\max\{k,7\}$.

\emph{Encoding.} $\Phi\ge0$ is an integer. $\Phi=0$ forces $S$ to increase
from $0$ to $1$ in steps $\zeta_l\in\{0,1\}$, so exactly one $\zeta_{l^*}=1$,
and then $(x_c,x_d)=(A_m,B_m)=(a_{l^*},b_{l^*})\in R_{cd}$. Thus $\Phi=0$
exactly at encodings of multicolored cliques, and each clique has exactly one
encoding because the pairs in $R_{cd}$ are distinct. At an encoding the
tie-break equals the weight of the clique's edge set. If a clique exists and
the minimum-weight clique is unique, $\Psi_w$ has a unique minimizer, value at
most $W_0-1$, and integral values; every point with $\Phi\ge1$ has
$\Psi_w\ge W_0$.

\emph{Conditioning.} All domains have length at most $N$, and there are
$n_v\le k+2k^2N^2$ variables, so $g\ge1/(n_vN^2)$; the diagonal Hessian entries
are at most $W_0(2+4N^2+2k)$. Hence $\kappa\le W_0(2+4N^2+2k)n_vN^2$, and both
$\kappa$ and $I$ are at most $(kN)^{c_2}$ for an absolute $c_2$.

\emph{Reduction.} Run the assumed algorithm for
$f(p)((kN)^{2c_2})^{p/\psi(p)}$ steps. If it outputs a feasible $\hat z$ with
$\Psi_w(\hat z)<W_0$, then $\Phi(\hat z)=0$ and $\hat z$ encodes a clique; answer
yes. Otherwise answer no. There are no false positives. If a clique exists,
Lemma~\ref{lem:isolation} (with $W=R$ and $\mathcal F$ the edge sets of
cliques) makes the minimum-weight clique unique with probability at least
$1/2$, and then the algorithm halts in time with $\Psi_w(\hat z)\le W_0-1/2$.
This is a one-sided-error algorithm for multicolored clique with running time
$f'(k)N^{o(k)}$, which rETH excludes \cite{ChenEtAl2006} (the reductions
behind that bound are deterministic, so they also exclude randomized
algorithms under rETH).
\end{proof}

